\documentclass{ctexart}
\usepackage{xeCJK}
\usepackage{xpinyin}
\usepackage{amsmath,amssymb,graphicx,booktabs,hyperref}
\usepackage{url}

\title{中文排版边界测试}

\author{测试作者}

\begin{document}

\maketitle

\begin{abstrct}
这是 Corpus \#23 中文排版边界测试的完整内容。
焦点：xeCJK、拼音、繁体字、生僻字。
\end{abstrct}

\section{引言}

本文测试中文排版的各种边界情况。

\subsection{简体中文内容}

简体中文内容包含常用汉字和标点符号。标点挤压测试：这是中文引号测试"双引号"，'单引号'。

\subsection{繁体字内容}

繁體字測試：這是繁體中文的內容，包含常用繁體字形。

\subsection{生僻字测试}

生僻字：靐（三个雷）、龘（三条龙）、齉（鼻子不通）、䶮（龙的一种）等 Unicode 生僻汉字。

\section{数学公式中的中文}

数学公式 $E = mc^2$ 与中文混合排版：这是中文标点与数学公式 $a^2 + b^2 = c^2$ 的混合测试。

\section{表格中的中文}

\begin{table}[ht]
\centering
\caption{中文表格示例}
\label{tab:chinese}
\begin{tabular}{lcc}
  \toprule
  序号 & 姓名 & 备注 \\
  \midrule
  1 & 张三 & 高级工程师 \\
  2 & 李四 & 产品经理 \\
  3 & 王五 & 技术总监 \\
  \bottomrule
\end{tabular}
\end{table}

\section{结论}

中文排版测试完成。

\bibliographystyle{gbt7714-numerical}
\bibliography{refs}

\end{document}
